%% Example CV - Angelo Marino Carmollingo
%% Comprehensive overview of all competencies, experience, and achievements.
%% Use as a master reference when tailoring role-specific CVs.
%%
%% SETUP: Run /setup to populate this with your actual information.
%% Compile with: lualatex main_example.tex
%% (pdflatex often fails on modern MiKTeX with fontawesome5 font-expansion errors.)

\documentclass[11pt,a4paper,sans]{moderncv}
\moderncvstyle{banking}
\moderncvcolor{blue}

% Force the name and section headings to render in moderncv blue (color1).
% Default banking leaves them black: moderncvstylebanking.sty's \colorlet
% copies (not aliases) the pre-scheme accent colour, so the name colours are
% frozen before \moderncvcolor runs. Re-let them after. \namefont is the hook
% every name-style macro routes through, so this also works on moderncv 2.3.1
% (Debian/Ubuntu apt), which has no \firstnamestyle/\lastnamestyle at all.
\renewcommand*{\namefont}{\fontsize{34}{36}\bfseries\upshape}
\colorlet{firstnamecolor}{color1}
\colorlet{lastnamecolor}{color1}
\colorlet{namecolor}{color1}
\renewcommand*{\sectionstyle}[1]{{\sectionfont\color{color1}#1}}

\usepackage[utf8]{inputenc}
% pdflatex fallback only (the documented engine is lualatex, which skips this
% branch). Without T1 font encoding pdflatex builds accented letters with
% \accent, and the PDF text layer stores them decomposed - `e` + U+0300 rather
% than U+00E8 - so an ATS keyword match on "Genève" fails while the page looks
% right. moderncv 2.5 loads T1 itself under pdflatex; 2.3.1 (Debian/Ubuntu apt)
% does not. \ifpdftex comes from iftex, which every moderncv version loads.
\ifpdftex\usepackage[T1]{fontenc}\fi
% moderncv loads hyperref itself in an \AtEndPreamble hook, so \hypersetup
% must go in an \AtEndPreamble of our own: on moderncv < 2.4 a top-level
% \usepackage{hyperref} clashes with the class's own
% \RequirePackage[unicode]{hyperref}. From 2.4.0 the class passes its options
% through \PassOptionsToPackage instead, which is what removes that clash.
\AtEndPreamble{\hypersetup{
    colorlinks=true,
    linkcolor=blue,
    filecolor=magenta,
    urlcolor=blue,
    pdftitle={Angelo Marino Carmollingo - CV},
    % Keep pdfpagemode=UseNone: this block runs after moderncv's own
    % \AtEndPreamble (moderncv.cls sets pdfpagemode there), so a FullScreen
    % value here would win and open every CV in fullscreen presentation mode.
    pdfpagemode=UseNone,
}}
\usepackage[scale=0.80]{geometry}
\usepackage{import}
\usepackage{needspace}  % evita titoli \cventry orfani a fine pagina

% personal data
\name{Angelo Marino}{Carmollingo}
% If you have no address to list, DELETE this whole line. \address{}{}{} fails
% with "There's no line here to end" on every moderncv version.
\address{Fano (PU), Italy}{}{}
\phone[mobile]{+39 375 534 4587}
\email{angelomarino.carmollingo@ingpec.eu}
\extrainfo{\href{https://linkedin.com/in/angelo-marino-carmollingo}{LinkedIn}, \href{https://github.com/AMCSoft}{GitHub}}

\begin{document}

\makecvtitle

% ============================================================
%     PROFILE STATEMENT
% ============================================================

\vspace{6pt}
\small{Software engineer with five years in domains where performance and correctness are not negotiable: flight-grade ISO-C99 for aerospace, embedded Linux for vehicle-to-vehicle communication, and AUTOSAR automotive software. Fluent in C and C++ at the level where allocation, layout and lifetimes are a default concern. Comfortable across the stack from performance-critical code to on-hardware validation. MSc in Computer Engineering, Politecnico di Torino.}
% NOTE: /apply replaces this statement per posting - see 05-cv-templates.md for the four variants.

% ============================================================
%     CORE COMPETENCIES
% ============================================================

\section{Core Competencies}
\vspace{1pt}
\begin{itemize}

\item \textbf{Languages}: C (ISO-C99, embedded, flight-grade), C++ (C++14/C++17, STL), Python, Bash, SQL.

\item \textbf{Embedded \& Systems}: Embedded Linux, Linux kernel manipulation, systemd modules, memory management in constrained environments, Bluetooth and LAN stacks.

\item \textbf{Automotive}: AUTOSAR (RTE, BSW, SWC), CAN communication, DaVinci Configurator and DaVinci Developer, bench testing on real hardware.

\item \textbf{Aerospace \& Safety-Critical}: PUS (Packet Utilization Standard), non-volatile memory filesystems, code documentation and validation to aerospace standards, CanTaTa, OsvALdo.

\item \textbf{Tooling \& Process}: Git, GitLab, GitHub, Jira, automated testing (CppUnitTest, Cantata, Python test harnesses), Qt/QML, OpenGL. Agile/SCRUM as practitioner and as Scrum Master.

\end{itemize}

% ============================================================
%     PROFESSIONAL EXPERIENCE
% ============================================================

\section{Professional Experience}
\vspace{3pt}
\begin{itemize}

% --- Most Recent Role ---
\needspace{5\baselineskip}
\item{\cventry{2024-Present}{Freelance Software Engineer}{Self-employed (main client: KooSys GmbH)}{Fano, Italy}{}{\vspace{1pt}
\begin{itemize}
    \item {AUTOSAR environment set-up and CAN communication for automotive projects}
    \item {RTE configuration with the DaVinci Configurator and DaVinci Developer toolchain}
    \item {Embedded software development in C, with Python scripting driving automated tests of the C codebase}
    \item {\textbf{Scrum Master} for the delivery team for approximately two years, running ceremonies, sprint planning and estimation while remaining a hands-on engineer}
    \item {\textbf{Project Manager / Technical Lead} since 2026: technical leadership and project ownership on top of the engineering role}
    \item {Delivery in Agile SCRUM using GitLab and Jira, remote from Italy}
\end{itemize}}}

\vspace{3pt}

% --- Previous Role ---
\needspace{5\baselineskip}
\item{\cventry{2022-2023}{Software Engineer}{DXC Luxoft}{Turin, Italy \& Berlin, Germany}{}{\vspace{1pt}
\begin{itemize}
    \item {Embedded Linux for vehicle-to-vehicle (V2V) communication between trucks: kernel manipulation, systemd modules, Bluetooth and LAN functionality}
    \item {C++14/C++17 and C99 development, with Python and Bash tooling}
    \item {Validation on real hardware benches; Agile SCRUM delivery}
\end{itemize}}}

\vspace{3pt}

% --- KooSys (employee, pre-freelance) ---
\needspace{5\baselineskip}
\item{\cventry{2023-2024}{Software Engineer}{KooSys GmbH}{Regensburg, Germany}{}{\vspace{1pt}
\begin{itemize}
    \item {RTE, BSW and SWC configuration and coding for next-generation automotive software}
    \item {DaVinci Toolchain; bench testing on real hardware}
\end{itemize}}}

\vspace{3pt}

% --- Earlier Role ---
\needspace{5\baselineskip}
\item{\cventry{2021-2022}{Junior Software Engineer}{Intelligentia}{Benevento, Italy}{}{\vspace{1pt}
\begin{itemize}
    \item {Developed PUS (Packet Utilization Standard) library code in ISO-C99 for defense and aerospace, documented and validated to flight-grade standards}
    \item {Designed and implemented, solo, a complete filesystem for non-volatile memory in a safety-critical spacecraft context (master's thesis)}
\end{itemize}}}

\vspace{3pt}

% --- Soft.Lab. (trainee) ---
\needspace{5\baselineskip}
\item{\cventry{2020-2021}{Trainee Software Engineer}{Soft.Lab. Srl}{Ponte (BN), Italy}{}{\vspace{1pt}
\begin{itemize}
    \item {Geomatic calculation and simulation software in C++14 with the STL and OpenGL; 3D scenarios}
    \item {Desktop applications with Qt and QML; automated testing with CppUnitTest}
\end{itemize}}}

\end{itemize}

% ============================================================
%     EDUCATION
% ============================================================

\section{Education}
\vspace{1pt}
\begin{itemize}

\needspace{5\baselineskip}
\item{\cventry{2020-2022}{MSc in Computer Engineering}{Politecnico di Torino}{Turin, Italy}{}{\vspace{1pt}
Software Development track. Thesis: a complete filesystem for non-volatile memory applications in a safety-critical (spacecraft) environment, designed and implemented solo at Intelligentia.
}}

\vspace{3pt}

\needspace{5\baselineskip}
\item{\cventry{2017-2020}{BSc in Computer Engineering}{Politecnico di Torino}{Turin, Italy}{}{\vspace{1pt}
Computer engineering fundamentals: algorithms and data structures, operating systems, computer architecture.
}}

\end{itemize}

% ============================================================
%     LANGUAGES
% ============================================================

\section{Languages}
\vspace{1pt}
\begin{itemize}
\item {Italian (native), English (professional working proficiency).}
\end{itemize}

% ============================================================
%     PUBLICATIONS (optional)
% ============================================================

% No publications on record - section commented out so no placeholder can reach a compiled PDF.
% Uncomment and fill if this changes.
% \section{Publications}
% \vspace{1pt}
% \begin{itemize}
% \item {Author(s) (Year). Title. Journal/Conference.}
% \end{itemize}

% ============================================================
%     HONORS AND AWARDS (optional)
% ============================================================

% No awards on record - section commented out for the same reason.
% \section{Honors and Awards}
% \vspace{1pt}
% \begin{itemize}
% \item{\textbf{Award Name} -- Event/Organization (Year).}
% \end{itemize}

% ============================================================
%     REFERENCES
% ============================================================

\section{References}
\vspace{1pt}
\begin{itemize}
\item{Available upon request.}
\end{itemize}

\end{document}
